\begin{lemma}[MetaCIC parallel-diamond canonical closurestatus reader]
\label{lem:metacic-parallel-diamond-frontier-canonical-closurestatus-reader}
The canonical paper status for $\MetaCICParallelDiamondFrontierUp$ is the bridge-checked closurestatus block in \texttt{closurestatus\_current.tex}. Any seed or top-transition reader that enters through the hub must discharge its status read through the child obligation surface
$
\mathsf{BEDC.Derived.MetaCICParallelDiamondFrontierUp.MetacicParallelDiamondFrontierObligations}
$
and then through that closurestatus block.
\end{lemma}
\begin{proof}
The hub displays the seed, bridged, and mature route markers, but its body is only a routing list. Therefore a status consumer must leave the hub and read the child files that carry the frontier evidence.

First read the checked obligation surface recorded in \autoref{thm:metacic-parallel-diamond-frontier-obligations}. It exposes the retained residual-substitution premise, critical-pair row, candidate-join row, residual-frontier row, closed-checker row, closed-normal fragment row, bounded-check row, obstruction row, transport, replay, provenance, and local naming before any status marker is consumed.

Then read the closurestatus child. Its block records $\matureClosure$, $\bridgeCheckedV$, and \texttt{bridgeChecked} for the same Lean obligation target, so the seed and top-transition routes point to that canonical child record rather than to a second hub-level status body.
\end{proof}
